\chapter{Literature Review}
\label{ch:literature}

\section{Introduction}
This chapter surveys the work on which the \projname\ builds. The survey is
organised around the two halves of the problem the system addresses. The first
half is scientific: what the Air Quality Index measures, how air quality has been
forecast, and which modelling families are appropriate for a tabular,
weather-driven, multi-horizon regression problem. The second half is
operational: how a trained model is kept alive in production, which is the
concern of the literature on feature stores, machine-learning operations,
distribution drift and model governance. A third, shorter strand covers
explainability and the honest evaluation of forecasts, both of which the system
treats as first-class requirements rather than afterthoughts. The chapter closes
with a literature matrix, an analysis of the gap this project addresses, and a
statement of how each body of work informed the design.

\section{Air Quality, Health and the Air Quality Index}

\subsection{The Health Burden}
The World Health Organization's global air-quality guidelines~\cite{ref:who}
substantially tightened the recommended limits for fine particulate matter in
2021, setting an annual mean guideline of 5\,$\mu$g/m$^3$ for
PM\textsubscript{2.5} and a 24-hour guideline of 15\,$\mu$g/m$^3$, on the basis
of accumulated evidence linking exposure to cardiovascular disease, respiratory
disease, adverse birth outcomes and premature mortality. The State of Global Air
report~\cite{ref:sga} places air pollution among the leading risk factors for
death worldwide and identifies South Asia as the region carrying the heaviest
exposure. Annual world air-quality reporting~\cite{ref:iqair} has for several
consecutive years ranked Pakistan among the most polluted countries measured,
with several Punjabi cities appearing among the most polluted urban areas
globally. A World Bank assessment of the economic cost of outdoor air pollution
in Pakistan~\cite{ref:worldbank} quantified the burden in terms of health
expenditure and lost productivity and argued for monitoring and forecasting
capacity as a prerequisite for any policy response.

Two features of this burden are directly relevant to the design of a forecasting
system. The first is that it is dominated by particulate matter rather than by
gaseous pollutants, which justifies the decision, described in
Chapter~\ref{ch:impl}, to compute the index from PM\textsubscript{2.5} and
PM\textsubscript{10} where the index has to be computed at all. The second is
that it is strongly seasonal: the exposure is concentrated in a winter smog
season driven by temperature inversions, crop residue burning and reduced mixing
height. Seasonality of that strength is exactly what a model with calendar and
weather features can learn, and exactly what makes a static model decay.

\subsection{The Air Quality Index}
The Air Quality Index used in this project is the United States Environmental
Protection Agency index, defined in the agency's technical assistance document
for daily air-quality reporting~\cite{ref:epa}. The index is not measured; it is
computed. For a pollutant concentration $C$ falling in the breakpoint interval
$[C_{lo}, C_{hi}]$, whose corresponding index interval is $[I_{lo}, I_{hi}]$, the
sub-index is the piecewise-linear interpolation

\begin{equation}
I = \frac{I_{hi} - I_{lo}}{C_{hi} - C_{lo}}\,(C - C_{lo}) + I_{lo},
\end{equation}

and the reported index is the \emph{maximum} of the pollutant sub-indices, on the
principle that health risk is governed by the worst offender rather than by an
average. The index is divided into six named health categories, from
\emph{Good} to \emph{Hazardous}, each carrying standard advisory language. This
categorical structure matters for the interface: a user does not need to
interpret the number 163, but does need to know that it means ``unhealthy'' and
that outdoor exertion should be limited.

An important practical subtlety, and a common source of error in student
implementations, is that the EPA defines gaseous breakpoints in parts per billion
or parts per million while most data services report gas concentrations in
micrograms per cubic metre. Applying the tables without a molar-mass conversion
produces a wildly inflated sub-index. This project's treatment of that issue is
described in Section~\ref{sec:impl-aqi}.

\section{Approaches to Air-Quality Forecasting}

\subsection{Deterministic Chemical Transport Modelling}
The physically grounded approach simulates emission, advection, chemical
transformation and deposition over a gridded domain. The Copernicus Atmosphere
Monitoring Service~\cite{ref:cams} operates such a system globally and publishes
composition forecasts, and regional implementations built on coupled
meteorology-chemistry models are the basis of most national air-quality
services. These models are indispensable for scientific attribution because they
represent the physics explicitly. They are, however, impractical for a project of
this kind: they require a maintained emissions inventory, meteorological boundary
conditions and substantial computing capacity, none of which is available for
free. The present work consumes the \emph{output} of such systems indirectly,
since the reanalysis and forecast products it ingests are themselves derived from
them.

\subsection{Statistical Time-Series Modelling}
The classical alternative fits an autoregressive integrated moving-average model,
or a seasonal variant, to the pollutant series at a single site. Hyndman and
Athanasopoulos~\cite{ref:fpp} give the standard modern treatment. Such models are
cheap, interpretable and often surprisingly hard to beat at very short lead
times, because pollutant series are strongly autocorrelated. Their limitation is
structural: a univariate model extrapolates the series' own history and cannot
consume a weather forecast, so its skill decays towards the series mean as the
horizon lengthens. This is precisely the regime in which the present system must
operate, since a three-day forecast is dominated by the weather that has not yet
happened rather than by the pollution that already has.

The same literature supplies the baseline against which any forecaster must be
judged. The naive or \emph{persistence} forecast - predicting that the value at
$t+h$ equals the value at $t$ - is the standard reference point, and Hyndman and
Athanasopoulos are emphatic that a model failing to beat it has demonstrated
nothing. This project therefore scores persistence explicitly in every training
run and at every horizon.

\subsection{Machine-Learning Regression}
The dominant contemporary approach treats forecasting as supervised regression
from engineered features to a future concentration. Yu et al.~\cite{ref:raq}
demonstrated a random-forest approach to urban air-quality prediction from
heterogeneous sensing data, and Pan~\cite{ref:xgbpm} reported gradient-boosted
trees performing strongly on hourly PM\textsubscript{2.5} prediction. The
appeal of the approach for this problem is that the strongest predictors of
future pollution - future wind, temperature, humidity and precipitation - are
themselves available as a free forecast, so the model can be given genuine
information about the period it is predicting rather than being restricted to
extrapolating the past.

The two tree ensembles used here rest on well-established foundations. Random
forests~\cite{ref:rf} reduce variance by averaging decorrelated trees grown on
bootstrap samples with randomised feature subsets. Gradient boosting~\cite{ref:gb}
instead fits trees sequentially to the residuals of the current ensemble,
performing gradient descent in function space; XGBoost~\cite{ref:xgboost} is the
scalable, regularised implementation of that idea that has become the default
choice for tabular problems, and it is the champion model in this system. Ridge
regression~\cite{ref:ridge} is included as a regularised linear reference point,
which is useful diagnostically: a large gap between the linear model and the tree
ensembles is evidence that the relationship being learned is genuinely non-linear
and interactive.

\subsection{Deep Sequence Models}
Recurrent networks, and in particular the long short-term memory
architecture~\cite{ref:lstm}, consume a sequence directly and can in principle
learn temporal structure that a tabular model must be given explicitly through
lag features. Li et al.~\cite{ref:lstmpm} applied an LSTM to air-pollutant
concentration prediction and reported improvements over shallower baselines, and
a substantial literature has followed in that direction.

Against this, a body of careful comparative work argues that deep architectures do
not dominate on tabular problems. Shwartz-Ziv and Armon~\cite{ref:tabnotall}
found that gradient-boosted trees outperformed several purpose-built deep tabular
architectures across a range of datasets, and Grinsztajn et
al.~\cite{ref:whytrees} analysed \emph{why}: tree ensembles cope better with
irregular target functions, are less affected by uninformative features, and are
not biased towards rotationally invariant solutions in a space where the
individual features carry meaning. This project treats the question empirically
rather than by assumption: an LSTM was implemented and evaluated on the same task
and the same chronological split as the tree ensembles, and the resulting
comparison, reported in Chapter~\ref{ch:testing}, informed the decision about
which model to promote.

\subsection{Multi-Horizon Forecasting Strategies}
When a forecast is required at many lead times, three strategies are available:
\emph{recursive}, in which a one-step model is applied repeatedly to its own
output; \emph{direct}, in which a separate model is trained for each horizon; and
the hybrid or multi-output variants. Ben Taieb et al.~\cite{ref:multistep}
compared these strategies systematically and identified the essential trade-off:
the recursive strategy accumulates its own errors as the horizon grows, while the
direct strategy avoids that accumulation at the cost of training and maintaining
one model per horizon.

The design adopted here is a variant of the direct strategy that avoids its cost.
The horizon $h$ is supplied to a single model as an input feature, so one global
estimator serves every lead time. This eliminates error accumulation, requires
exactly one artefact to be versioned, registered, explained and monitored, and
allows every horizon to be trained on the whole dataset rather than on a slice of
it. The design and its consequences are developed in Section~\ref{sec:design-mh}.

\section{Honest Evaluation of Forecasts}
Two methodological points from the forecasting literature shape how this system
reports its accuracy.

The first concerns validation. Tashman~\cite{ref:tashman} reviewed out-of-sample
testing practice and argued for rolling-origin evaluation over a single holdout,
and Bergmeir and Benítez~\cite{ref:bergmeir} examined the use of cross-validation
for time-series predictors and set out the conditions under which random $k$-fold
partitioning is invalid. The essential point is that a randomly chosen validation
set for a time series lets the model train on the future and test on the past,
which flatters it. This project therefore splits strictly chronologically and,
beyond that, reports a five-fold walk-forward backtest in which each fold trains
on the past and tests on the window that follows.

The second concerns reporting granularity. A multi-horizon forecaster summarised
by one averaged error figure conceals the fact that skill decays sharply with
lead time; a single number that mixes a one-hour and a 72-hour forecast is
dominated by whichever horizon happens to be more numerous in the validation set,
and it overstates usefulness at the horizons a user actually cares about. This
project therefore reports RMSE, MAE and $R^2$ separately at every modelled
horizon, alongside the persistence baseline at that same horizon.

\section{Uncertainty Quantification}
A point forecast without an interval invites false confidence. The most
theoretically satisfying framework for attaching a distribution-free interval to
an arbitrary regressor is conformal prediction, introduced by Vovk, Gammerman and
Shafer~\cite{ref:conformal}; the computationally practical inductive or
\emph{split} variant is due to Papadopoulos et al.~\cite{ref:inductive}, and Lei
et al.~\cite{ref:leiconformal} give a modern treatment with finite-sample
coverage guarantees. In its simplest form, split conformal prediction fits the
model on a training set, computes residuals on a disjoint calibration set, and
constructs the interval from the empirical quantiles of those residuals; the
resulting interval attains the nominal coverage under exchangeability.

The intervals in this system are constructed in exactly this spirit: residual
quantiles are computed on the held-out chronological validation set and stored
per horizon, so that the interval widens with lead time as the model's actual
errors do. The exchangeability assumption is, strictly, not satisfied for a
non-stationary time series, and Chapter~\ref{ch:testing} discusses the
consequences of that honestly rather than claiming a guarantee the construction
does not earn.

\section{Explainable Machine Learning}
Ribeiro et al.~\cite{ref:lime} introduced LIME, which explains an individual
prediction by fitting an interpretable surrogate in its neighbourhood. Lundberg
and Lee~\cite{ref:shap} subsequently unified several attribution methods under
the framework of Shapley values from cooperative game theory, producing SHAP: the
unique additive attribution satisfying local accuracy, missingness and
consistency. The practical obstacle to Shapley values is their exponential cost,
which Lundberg et al.~\cite{ref:treeshap} removed for tree ensembles with
TreeSHAP, a polynomial-time exact algorithm that also supports aggregation of
local attributions into global importance.

Two properties make SHAP the right choice here. It is additive, so the
attributions for a single forecast sum exactly to the difference between that
forecast and the model's base value - which means an explanation can be rendered
as an arithmetic statement a non-specialist can follow, rather than as a
ranking whose units are unclear. And TreeSHAP is exact and fast for the
gradient-boosted champion, so an explanation can be computed inside a batch job
without materially lengthening it. Section~\ref{sec:impl-shap} describes how the
attributions are translated into natural language.

\section{Feature Stores and the Feature-Training-Inference Pattern}
Sculley et al.~\cite{ref:techdebt} made the observation that frames the
operational half of this project: in a real machine-learning system the model is
a small fraction of the code, and the surrounding infrastructure - data
collection, feature extraction, configuration, serving, monitoring - is where
the complexity and the accumulated debt actually live. They identified
\emph{training-serving skew}, in which features are computed one way for
training and another way for inference, as a recurring and particularly damaging
failure mode.

The feature store emerged as the industrial answer to that problem. Introduced
publicly as a component of Uber's Michelangelo platform~\cite{ref:michelangelo},
it is a versioned, queryable repository of engineered features with a defined
schema, primary keys and event time, written by feature pipelines and read by
both training and inference. Because both sides read the same table, skew is
eliminated by construction rather than by discipline. Hopsworks, the platform used
in this project, provides such a store together with an integrated model
registry~\cite{ref:hopsworks}.

The architectural pattern that follows from this is the
Feature-Training-Inference decomposition~\cite{ref:fti}, in which a machine
learning system is expressed as exactly three pipelines that share no code path
and communicate only through the feature store and the model registry. The
pattern's value for this project is specific and practical: because the only
durable state lives in managed services, the compute can be entirely ephemeral,
which is what allows the whole system to run on a free scheduler.

\section{Machine-Learning Operations}
Breck et al.~\cite{ref:mltestscore} proposed a rubric for production readiness in
machine-learning systems, scoring tests across four areas - features and data,
model development, infrastructure, and monitoring - and observing that most
deployed systems score poorly on the last two. Their rubric supplies, in effect,
a checklist against which the present system can be assessed, and several of its
items map directly onto components implemented here: that the model is tested
against a simpler baseline, that training is reproducible and scheduled, that
model quality is validated before serving, and that the model can be rolled back.

Two of those items deserve individual attention.

\subsection{Gated Promotion and the Champion-Challenger Pattern}
The champion-challenger pattern originates in credit scoring and marketing
analytics, where it is standard practice that an incumbent model continues to
make decisions until a challenger has demonstrated superiority on an agreed
metric; Siddiqi~\cite{ref:siddiqi} documents it in that setting alongside the
monitoring practice discussed below. Transplanted into an automated retraining
pipeline, it answers a question that automation creates: if a system retrains
every night, what stops it from deploying whatever it happened to produce last
night? Without a gate, a retrain on an unrepresentative window can quietly
degrade a live service. The gate implemented here, described in
Section~\ref{sec:impl-gate}, promotes a challenger only if it improves on the
champion's validation RMSE, and records every decision - including refusals -
in a persistent leaderboard.

\subsection{Drift Detection}
Gama et al.~\cite{ref:driftsurvey} survey concept drift and its adaptation
strategies, distinguishing between change in the input distribution and change in
the relationship between inputs and target, and between abrupt, gradual and
recurring change. For a seasonal environmental system, \emph{recurring} drift is
the normal condition rather than a fault, which is an important interpretive
point: a drift alarm in November does not mean the system is broken.

The measure used in this project is the Population Stability Index, long
established in credit-risk monitoring~\cite{ref:siddiqi}. For a reference
distribution binned into $B$ bins with proportions $r_b$, and a current
distribution with proportions $c_b$ over the same bin edges,

\begin{equation}
\mathrm{PSI} = \sum_{b=1}^{B} (c_b - r_b)\,\ln\!\left(\frac{c_b}{r_b}\right),
\end{equation}

with the conventional interpretation that values below 0.10 indicate a stable
population, values between 0.10 and 0.25 moderate shift, and values above 0.25
significant shift. PSI is attractive here because it is non-parametric, cheap,
computed per feature so that the responsible variable is identified, and
interpretable against a threshold that does not need to be tuned. Its weakness -
that it measures input drift only, and says nothing about whether accuracy has
actually suffered - is why this system pairs it with direct measurement of
realised forecast error against observed outcomes.

\section{Language Models as an Interface Layer}
The final component of the system exposes its forecasts to a large language
model, so that a user can ask a question in ordinary language rather than reading
a chart. The relevant methodological concern is grounding: a model asked about
air quality will otherwise produce a plausible number from its parameters rather
than the system's actual forecast. The mitigation adopted is tool use, in which
the model is given a declared set of functions and returns a structured call
rather than prose, and the answer is composed from the function's real return
value. The Model Context Protocol~\cite{ref:mcp} standardises the declaration and
invocation of such tools across clients, which allows the same four grounded
functions to serve both the dashboard's advisor and any external protocol client.

\section{Literature Review Matrix}
Table~\ref{tab:litmatrix} condenses the survey, relating each strand of work to
what it contributes and to what it leaves open for this project.

\begin{xltabular}{\textwidth}{P{2.9cm} P{3.4cm} Y}
\caption{Literature review matrix.}\label{tab:litmatrix}\\
\toprule
\bfseries Strand & \bfseries Representative work & \bfseries Contribution and limitation for this project \\
\midrule
\endfirsthead
\toprule
\bfseries Strand & \bfseries Representative work & \bfseries Contribution and limitation for this project \\
\midrule
\endhead
\bottomrule
\endlastfoot
Health burden and index definition & WHO~\cite{ref:who}; EPA~\cite{ref:epa};
IQAir~\cite{ref:iqair} & Defines what is being predicted and why it matters; establishes
that the burden is particulate-dominated and seasonal. Does not address forecasting. \\
Chemical transport modelling & CAMS~\cite{ref:cams} & Physically rigorous and the
scientific reference. Requires an emissions inventory and computing capacity that a
free-tier project cannot supply. \\
Statistical time series & Hyndman and Athanasopoulos~\cite{ref:fpp} & Supplies the
persistence baseline and the discipline of comparing against it. Univariate models cannot
consume a weather forecast, so they fade at long horizons. \\
Tree ensembles & Breiman~\cite{ref:rf}; Friedman~\cite{ref:gb};
Chen and Guestrin~\cite{ref:xgboost} & The modelling core: strong on tabular,
interactive, weather-driven data and cheap to retrain daily. Requires temporal
structure to be supplied explicitly through engineered features. \\
Deep sequence models & Hochreiter and Schmidhuber~\cite{ref:lstm};
Li et al.~\cite{ref:lstmpm} & Learns temporal structure directly and is a legitimate
contender. Costly to train, single-horizon as usually formulated, and harder to explain. \\
Trees versus deep learning on tabular data & Shwartz-Ziv and
Armon~\cite{ref:tabnotall}; Grinsztajn et al.~\cite{ref:whytrees} & Provides the
principled reason to test rather than assume, and predicts the outcome observed here. \\
Multi-horizon strategy & Ben Taieb et al.~\cite{ref:multistep} & Frames the direct versus
recursive choice. Leaves open the practical cost of one model per horizon, which the
horizon-as-feature design resolves. \\
Validation methodology & Tashman~\cite{ref:tashman};
Bergmeir and Benítez~\cite{ref:bergmeir} & Justifies chronological splitting and
walk-forward backtesting over random cross-validation. \\
Uncertainty & Vovk et al.~\cite{ref:conformal};
Lei et al.~\cite{ref:leiconformal} & Gives a distribution-free construction for
intervals. Its exchangeability assumption is imperfect for a seasonal series. \\
Explainability & Lundberg and Lee~\cite{ref:shap};
Lundberg et al.~\cite{ref:treeshap} & Exact, fast, additive attributions for tree models,
which can be rendered arithmetically in plain language. \\
Feature stores and FTI & Sculley et al.~\cite{ref:techdebt};
Michelangelo~\cite{ref:michelangelo}; Dowling~\cite{ref:fti} & Supplies the architecture
that eliminates training-serving skew and permits stateless compute. Assumes a managed
store is available, which the free tier constrains. \\
MLOps practice & Breck et al.~\cite{ref:mltestscore};
Siddiqi~\cite{ref:siddiqi} & Provides the readiness rubric, the promotion gate and the
drift measure. Neither is specific to environmental forecasting. \\
Drift & Gama et al.~\cite{ref:driftsurvey} & Distinguishes recurring seasonal drift from
genuine decay, which is essential to interpreting this system's own alarms. \\
Grounded language interfaces & MCP~\cite{ref:mcp} & Standardises tool declaration so a
language model answers from real forecasts rather than from its parameters. \\
\end{xltabular}

\section{Research Gap Analysis}
Read together, the literature leaves four gaps that this project addresses.

\textbf{Gap 1 - Coverage.} Published air-quality forecasting work for Pakistan is
overwhelmingly concerned with Lahore or Karachi, and models are typically fitted
per site. A single global model conditioned on location, trained across 22 cities
spanning every province and territory, is not the usual formulation, and it
carries a specific benefit: a city with a shorter or noisier record borrows
strength from the rest.

\textbf{Gap 2 - Operation rather than experiment.} The overwhelming majority of
the modelling literature reports an accuracy figure for a model trained once.
Scheduled retraining, gated promotion, a durable registry and self-monitoring are
described in the MLOps literature and in industrial platform papers, but are
rarely instantiated in published air-quality work. This project's contribution is
to assemble the two.

\textbf{Gap 3 - Reporting that reflects use.} Where multi-horizon forecasts are
published, accuracy is usually reported as a single averaged number. Reporting
per horizon, against a per-horizon baseline, and under walk-forward backtesting,
gives a materially different - and less flattering, therefore more useful -
picture of what a three-day forecast is actually worth.

\textbf{Gap 4 - Explanation and uncertainty as delivered features.} SHAP and
conformal intervals are both well established as techniques, but they are usually
evaluated in isolation. Here they are delivered: every forecast in the public
interface carries an interval and a sentence explaining what drove it, and the
what-if simulator lets a user perturb the drivers and watch the model respond.

\section{Relevance to This Project}
Each strand of the survey translated into a concrete design decision.

The health and index literature fixed the target: the EPA AQI, with its six
categories and their advisory language, computed from particulate matter where it
must be computed at all. The comparison of forecasting traditions ruled out
physical simulation on resource grounds and univariate statistics on horizon
grounds, and pointed to supervised regression with forecast weather as inputs.
The tabular-learning literature justified making XGBoost the primary candidate
while requiring that the sequence-model alternative be measured rather than
dismissed. The multi-horizon literature motivated the horizon-as-feature
formulation. The validation literature dictated chronological splitting, a
persistence baseline at every horizon and a walk-forward backtest. The conformal
literature supplied the interval construction. The explainability literature
supplied TreeSHAP and the additive property that makes plain-language rendering
possible. The feature-store and FTI literature supplied the architecture, and the
MLOps literature supplied the promotion gate, the registry and the drift measure
that make the system self-operating. The tool-use literature supplied the
grounding discipline for the advisor.

\section{Summary of Literature Review}
This chapter reviewed the health and regulatory basis of the Air Quality Index,
the three traditions of air-quality forecasting and the reasons for choosing
supervised machine learning among them, the model families relevant to a tabular
multi-horizon regression and the evidence on trees versus deep networks, the
methodology of honest forecast evaluation and uncertainty quantification, the
explainability techniques appropriate to tree ensembles, and the operational
literature on feature stores, the Feature-Training-Inference pattern, gated
promotion and drift. It identified four gaps - coverage, operation, reporting
granularity, and delivered explanation - which the remainder of this thesis
addresses. The next chapter states the problem formally.
